\documentclass[../数学分析培优讲义.tex]{subfiles}

\begin{document}

\setcounter{chapter}{3}
\pagestyle{fancy}

\chapter{一元函数微分学}


\section{导数}


\subsection{导数的定义及可微性验证}

\begin{example}
设 $f(x)=|x|$,则 $f(x)$ 在 $x\neq0$ 处可导,且
\[
    f'(x)=\dfrac{|x|}{x}.
\]
\end{example}
\exampleblank

\begin{example}
证明:$f(x)=|x|^3$ 在 $x=0$ 处的三阶导数 $f'''(0)$ 不存在.
\end{example}
\exampleblank

\begin{example}
设
\[
f(x)=
\begin{cases}
    x^2\sin\dfrac{\pi}{x}, & x<0,\\
    A, & x=0,\\
    ax^2+b, & x>0,
\end{cases}
\]
其中 $A,a,b$ 为常数,试问 $A,a,b$ 为何值时 $f(x)$ 在 $x=0$ 处可导,为什么?并求 $f'(0)$.
\end{example}
\exampleblank

\begin{example}
设函数 $f(x)$ 在闭区间 $\lbrack 0,1\rbrack$ 上四次连续可微,$f(0)=0,f'(0)=0$,证明函数
\[
F(x)=
\begin{cases}
    \dfrac{f(x)}{x^2}, & 0<x\leq1,\\
    \dfrac{f''(0)}{2}, & x=0,
\end{cases}
\]
在 $\lbrack 0,1\rbrack$ 上二次连续可微.
\end{example}
\exampleblank

\begin{example}
设 $f(x)$ 在 $x_0$ 的某邻域内有定义.
\begin{enumerate}[label=(\arabic*)]
    \item 若 $f(x)$ 在 $x_0$ 处可导,试证
    \[
        \lim\limits_{h\to0}
        \dfrac{f(x_0+h)-f(x_0-h)}{2h}=f'(x_0);
    \]
    \item 反之,若上式左端极限存在,是否能推出 $f'(x_0)$ 存在?若结论成立,请证明;不成立请给出反例.
\end{enumerate}
\end{example}
\exampleblank

\begin{example}[\markstar]
证明 Riemann 函数
\[
R(x)=
\begin{cases}
    \dfrac{1}{q}, & x=\dfrac{p}{q},\\
    0, & x=0,1\text{ 或无理数},
\end{cases}
\]
处处不可微.
\end{example}
\exampleblank

\begin{proposition}[\markstar]
设函数
\[
f(x)=
\begin{cases}
    x^\alpha\sin\dfrac{1}{x}, & x\neq0,\\
    0, & x=0,
\end{cases}
\]
则:
\begin{enumerate}[label=(\arabic*)]
    \item $\alpha>0$ 时,在点 $x=0$ 连续;
    \item $\alpha>1$ 时,在点 $x=0$ 处可导;
    \item $\alpha>2$ 时,在点 $x=0$ 处连续可导.
\end{enumerate}
\end{proposition}

\begin{example}
求 $f(x)=[x]\sin\pi x$ 的单侧导数,并讨论可微性.
\end{example}
\exampleblank


\newpage

\subsection{高阶导数}


\methodsection{(A) 利用已知结论}

\begin{theorem}[\markstar]
常用的 $n$ 阶导数公式为
\[
    (x^k)^{(n)}=k(k-1)\cdots(k-n+1)x^{k-n},
    \qquad n\leq k,
\]
\[
    (e^x)^{(n)}=e^x,
    \qquad
    (a^x)^{(n)}=a^x(\ln a)^n,
\]
\[
    (\ln x)^{(n)}=(-1)^{n-1}\dfrac{(n-1)!}{x^n},
\]
\[
    (\sin x)^{(n)}=\sin\left(x+\dfrac{n\pi}{2}\right),
    \qquad
    (\cos x)^{(n)}=\cos\left(x+\dfrac{n\pi}{2}\right).
\]
\end{theorem}

\begin{remark}
因此有了以上常见 $n$ 阶导数公式,对于一般不易直接求 $n$ 阶导数的函数,我们优先考虑转化成如上形式函数.
\end{remark}

\begin{example}
计算 $y^{(n)}$,设
\begin{enumerate}[label=(\arabic*)]
    \item $y=\sin^6x+\cos^6x$;
    \item $y=\sin ax\cdot\sin bx$;
    \item $y=\dfrac{x^2+x+1}{x^2-5x+6}$.
\end{enumerate}
\end{example}
\exampleblank


\newpage

\methodsection{(B) 利用 Leibniz 公式}

\begin{example}
设 $f_1(x),f_2(x)$ 在 $x_0$ 及其附近有定义,在 $x_0$ 处有直到 $n$ 阶导数,记
\[
    N(f)=\sum\limits_{k=0}^n\dfrac{1}{k!}|f^{(k)}(x_0)|.
\]
试证明:
\[
    N(f_1f_2)\leq N(f_1)N(f_2).
\]
\end{example}
\exampleblank

\begin{example}
试证:
\[
    (a^2+b^2)^{\frac{n}{2}}e^{ax}\sin(bx+n\varphi)
    =\sum\limits_{i=0}^n C_{n}^{i}a^{n-i}b^i e^{ax}
    \sin\left(bx+\dfrac{\pi i}{2}\right),
\]
其中
\[
    \varphi=\arctan\dfrac{b}{a}.
\]
\end{example}
\exampleblank


\newpage

\methodsection{(C) 利用数学归纳法}

当高阶导数不易求出时,可考虑先求出前几阶导数,观察得出规律,进而利用数学归纳法证明.

\begin{example}[\markstar]
证明:若函数
\[
f(x)=
\begin{cases}
    e^{-\frac{1}{x^2}}, & x\neq0,\\
    0, & x=0,
\end{cases}
\]
则在 $x=0$ 处,
\[
    f^{(n)}(0)=0,
    \qquad n=1,2,\ldots.
\]
\end{example}
\exampleblank

\begin{example}[推广的 Leibniz]
设 $u_1(x),u_2(x),\ldots,u_k(x)$ 都是 $x$ 的 $n$ 次可微函数,试证:
\[
    (u_1u_2\cdots u_k)^{(n)}
    =\sum\limits_{r_1+\cdots+r_k=n}
    \dfrac{n!}{r_1!r_2!\cdots r_k!}
    u_1^{(r_1)}\cdots u_k^{(r_k)}.
\]
\end{example}
\exampleblank

\begin{example}[\markstar]
$f(x)$ 有任意阶导数,$f'(x)=(f(x))^2$,求 $f^{(n)}(x)\ (n>2)$.
\end{example}
\exampleblank


\newpage

\methodsection{(D) 利用递推公式}

当更高阶导数可以用低阶导数表示且易于推广到一般情况时,可考虑利用递推关系求更高阶导数.

\begin{example}
设 $f(x)=(\arcsin x)^2$,求 $f^{(n)}(0)$.
\end{example}
\exampleblank

\begin{example}
证明 Legendre 多项式
\[
    p_n(x)=\dfrac{1}{2^n n!}\bigl[(x^2-1)^n\bigr]^{(n)},
    \qquad n=0,1,2,\ldots,
\]
满足方程
\[
    (1-x^2)p_n''(x)-2xp_n'(x)+n(n+1)p_n(x)=0.
\]
\end{example}
\exampleblank


\newpage

\methodsection{(E) 利用 Taylor 展开}

\begin{example}
求 $f(x)=\arctan x$ 在 $x=0$ 处的各阶导数.
\end{example}
\exampleblank

\begin{example}
设
\[
f(x)=
\begin{cases}
    \dfrac{\sin x}{x}, & x\neq0,\\
    1, & x=0,
\end{cases}
\]
求 $f^{(k)}(0)$.
\end{example}
\exampleblank


\newpage

\subsection{利用导数求和}

\begin{example}
试求下列极限:
\begin{enumerate}[label=(\arabic*)]
    \item $\displaystyle
    I=\lim\limits_{x\to0}\dfrac{1}{x}\sum\limits_{i=1}^k f\left(\dfrac{x}{i}\right)$,
    已知 $f'(0)$ 存在,且 $f(0)=0,k\in\mathbb{N}$;
    \item $\displaystyle
    I=\lim\limits_{n\to\infty}n\left[\sum\limits_{i=1}^k f\left(a+\dfrac{i}{n}\right)-kf(a)\right]$,
    已知 $f'(a)$ 存在;
    \item $\displaystyle
    I=\lim\limits_{n\to\infty}
    \dfrac{\left(a+\dfrac{1}{n}\right)^n\left(a+\dfrac{2}{n}\right)^n\cdots\left(a+\dfrac{k}{n}\right)^n}{a^{nk}}$,
    $k\in\mathbb{N}$.
\end{enumerate}
\end{example}
\exampleblank

\begin{example}
试求下列和式:
\begin{enumerate}[label=(\arabic*)]
    \item $\displaystyle I_n=\sum\limits_{k=0}^n ke^{kx}$;
    \item $\displaystyle \theta_n(x)=\sum\limits_{k=1}^n k^2x^{k-1}$;
    \item $\displaystyle I_n=\sum\limits_{k=1}^n C_{n}^{k}k^2$;
    \item $\displaystyle I_n=\sum\limits_{k=0}^{2n}(-1)^k C_{2n}^{k}k^n$;
    \item $\displaystyle S_n(x)=\sum\limits_{k=1}^n\dfrac{1}{2^k}\tan\left(\dfrac{x}{2^k}\right)$.
\end{enumerate}
\end{example}
\exampleblank

\begin{example}
设 $f_1(x),f_2(x),\ldots,f_n(x)$ 在 $x=x_0$ 处可导,且 $f_k(x_0)\neq0\ (k=1,2,\ldots,n)$,
\[
    F(x)=\prod\limits_{k=1}^n f_k(x),
\]
则
\[
    \dfrac{F'(x_0)}{F(x_0)}
    =\sum\limits_{k=1}^n\dfrac{f_k'(x_0)}{f_k(x_0)}.
\]
\end{example}
\exampleblank


\newpage

\subsection{利用导数证恒等式}

\begin{example}
证明:
\[
    2\arctan x+\arcsin\dfrac{2x}{1+x^2}=\pi,
    \qquad 1<x<\infty.
\]
\end{example}
\exampleblank

\begin{example}
设 $f(x),g(x)$ 在 $(-\infty,+\infty)$ 上均不是常数,可导且有
\[
\begin{cases}
    f(x+y)=f(x)f(y)-g(x)g(y),\\
    g(x+y)=f(x)g(y)+g(x)f(y),
\end{cases}
\]
其中 $x,y\in(-\infty,+\infty)$,若 $f'(0)=0$,则
\[
    f^2(x)+g^2(x)=1.
\]
\end{example}
\exampleblank

\begin{example}
设 $f(x),g(x)$ 在 $(a,b)$ 上均不是常数,若有
\[
    f(x)+g(x)\neq0,
    \qquad
    f(x)g'(x)-f'(x)g(x)=0,
    \qquad x\in(a,b),
\]
则 $g(x)\neq0\ (x\in(a,b))$ 且
\[
    \dfrac{f(x)}{g(x)}\equiv C\text{(常数)}.
\]
\end{example}
\exampleblank


\newpage

\subsection{导函数的性质}

导函数有两个重要基本定理---Darboux 定理和单侧导数极限定理,进而可得导数的两个重要特性:\textcircled{1} 介值性;\textcircled{2} 无第一类间断点.

\begin{theorem}[\markstar\ Darboux 定理]
设 $f$ 在区间 $I$ 上可微,则 $f'$ 具有介值性.
\end{theorem}

\begin{corollary}[\markstar]
单调的导函数 $f'$ 必连续.
\end{corollary}

\begin{theorem}[\markstar\ 单侧导数极限定理]
设在区间 $\lbrack a,b)$ 上有定义的函数 $f$ 在 $(a,b)$ 上可微,又在点 $a$ 右连续,若导函数 $f'(x)$ 在点 $a$ 存在右侧极限
\[
    f'(a+0)=A,
\]
则 $f$ 在点 $a$ 也一定存在右侧导数 $f'_+(a)$,且成立
\[
    f'_+(a)=f'(a+0)=A.
\]
即 $f'(x)$ 在点 $a$ 右连续.
\end{theorem}

\begin{corollary}[\markstar]
设 $f$ 在 $(a,b)$ 上可微,则导函数 $f'(x)$ 不存在第一类间断点.
\end{corollary}

\begin{example}
设函数 $f(x)$ 在区间 $(-\infty,+\infty)$ 上二次可微,且有界,试证存在点 $x_0\in(-\infty,+\infty)$,使得 $f'(x_0)=0$.
\end{example}
\exampleblank

\begin{example}
设函数 $f$ 在 $(0,1)$ 有界、可导,但 $\displaystyle\lim\limits_{x\to0^+}f(x)$ 不存在,试证存在数列 $x_n\to0^+\ (n\to\infty)$,使得
\[
    \lim\limits_{n\to+\infty}f'(x_n)=0.
\]
\end{example}
\exampleblank

\begin{example}
\begin{enumerate}[label=(\arabic*)]
    \item 设 $f(x)$ 在 $(a,b)$ 上可导,且有 $x_i\in(a,b),\lambda_i>0\ (i=1,2,\ldots,n)$,
    \[
        \sum\limits_{i=1}^n\lambda_i=1,
    \]
    则存在 $\xi\in(a,b)$,使得
    \[
        \sum\limits_{i=1}^n\lambda_i f'(x_i)=f'(\xi);
    \]
    \item 设 $f(x)$ 在 $(a,b)$ 上可导,且有 $x_i<y_i$,$x_i,y_i\in(a,b)\ (i=1,2,\ldots,n)$,则存在 $\xi\in(a,b)$,使得
    \[
        \sum\limits_{i=1}^n[f(y_i)-f(x_i)]
        =f'(\xi)\sum\limits_{i=1}^n(y_i-x_i).
    \]
\end{enumerate}
\end{example}
\exampleblank

\begin{example}
设 $f(x)$ 在 $\lbrack a,b\rbrack$ 上二次可导,若 $f''(x)\neq0\ (a\leq x\leq b)$,则对任意 $\xi\in(a,b)$,存在 $\xi'$ 满足 $a<\xi<\xi'\leq b$,使得
\[
    f'(\xi)=\dfrac{f(\xi')-f(a)}{\xi'-a}.
\]
\end{example}
\exampleblank


\newpage

\subsection{广义导数 *}

\begin{definition}
$\Delta_hf(x)=f(x+h)-f(x-h)$ 称为 $f(x)$ 在点 $x$ 处的对称差,如下极限(若存在)称为 $f(x)$ 在点 $x$ 处的广义导数或对称导数,记作
\[
    f^{[\prime]}(x)
    =\lim\limits_{h\to0^+}\dfrac{\Delta_hf(x)}{2h}
    =\lim\limits_{h\to0^+}\dfrac{f(x+h)-f(x-h)}{2h}.
\]
\end{definition}

\begin{theorem}[\markstar]
设 $f$ 在点 $x$ 处可导,则 $f^{[\prime]}$ 存在且
\[
    f^{[\prime]}(x)=f'(x).
\]
\end{theorem}

\begin{remark}
该定理逆命题不成立,例如
\[
f(x)=
\begin{cases}
    x\sin\dfrac{1}{x}, & x\neq0,\\
    0, & x=0,
\end{cases}
\]
或 $f(x)=|x|$.
\end{remark}

\begin{definition}
如下的极限存在,称为 $f(x)$ 在点 $x$ 处有二阶广义导数(或称二阶对称导数),记作
\[
\begin{aligned}
    f^{[\prime\prime]}(x)
    &=f^{[2]}(x)\\
    &\overset{\triangle}{=}
    \lim\limits_{h\to0^+}\dfrac{\Delta_h\Delta_hf(x)}{(2h)^2}\\
    &=\lim\limits_{h\to0^+}
    \dfrac{\Delta_h[f(x+h)-f(x-h)]}{(2h)^2}\\
    &=\lim\limits_{h\to0^+}
    \dfrac{f(x+2h)-2f(x)-f(x-2h)}{4h^2}.
\end{aligned}
\]
\end{definition}

\begin{theorem}[\markstar]
在某点 $x$,若 $f(x)$ 有二阶导数,则二阶广义导数也存在,且值相等,即
\[
    f''(x)=f^{[2]}(x).
\]
\end{theorem}

\begin{remark}
逆命题同样不成立,例如
\[
f(x)=
\begin{cases}
    x^2\sin\dfrac{1}{x}, & x\neq0,\\
    0, & x=0,
\end{cases}
\]
或 $f(x)=x|x|$.
\end{remark}

\begin{example}[\markstar\ Schwarz 定理]
若 $f(x)$ 在 $\lbrack a,b\rbrack$ 上连续,$f(x)$ 的广义二阶导数 $f^{[2]}(x)$ 存在且恒为零,则 $f(x)$ 是一次线性函数,即 $f(x)=ax+b$.
\end{example}
\exampleblank

\begin{example}
若 $f(x)$ 在 $\lbrack a,b\rbrack$ 上连续,$f(a)<f(b)$,又设对一切 $x\in(a,b)$,
\[
    \lim\limits_{t\to0}\dfrac{f(x+t)-f(x-t)}{t}
\]
存在,用 $g(x)$ 表示这一极限值,则存在 $c\in(a,b)$,使得 $g(c)\geq0$.
\end{example}
\exampleblank


\newpage

\section{微分中值定理}


\subsection{构造辅助函数(中值问题)}

\noindent\textbf{(\markstar)} 看到 $f'(\xi)+f(\xi)g(\xi)$,应想到其原函数为
$\displaystyle f(x)e^{\int g(x)\,\mathrm{d}x}$,这是因为
\[
    \left[f(x)e^{\int g(x)\,\mathrm{d}x}\right]'
    =[f'(x)+f(x)g(x)]e^{\int g(x)\,\mathrm{d}x}.
\]
具体地,有如下常见体现:
\[
    mf(\xi)+nf'(\xi)
    \longleftrightarrow
    f(x)e^{\frac{m}{n}x}.
\]
\[
    mf(\xi)+n\xi f'(\xi)
    \longleftrightarrow
    x^m f^n(x).
\]
\[
    mf(\xi)-n\xi f'(\xi)
    \longleftrightarrow
    \dfrac{f^n(x)}{x^m}
    \quad\text{或}\quad
    \dfrac{x^m}{f^n(x)}.
\]
\[
    nf'(\xi)f(1-\xi)-mf(\xi)f'(1-\xi)
    \longleftrightarrow
    f^n(x)f^m(1-x).
\]
\[
    mf'(\xi)g(\xi)+nf(\xi)g'(\xi)
    \longleftrightarrow
    f^m(x)g^n(x).
\]
\[
    mf'(\xi)g(\xi)-nf(\xi)g'(\xi)
    \longleftrightarrow
    \dfrac{f^m(x)}{g^n(x)}.
\]
\[
    f(\xi)g''(\xi)-g(\xi)f''(\xi)
    \longleftrightarrow
    f'(x)g(x)-f(x)g'(x).
\]

\begin{example}
设 $f\in C(\lbrack a,b\rbrack)$,$g\in C(\lbrack a,b\rbrack)$,且均在 $(a,b)$ 上可导,若 $f(a)=f(b)=0$,则存在 $\xi\in(a,b)$,使得
\[
    f'(\xi)+g'(\xi)f(\xi)=0.
\]
\end{example}
\exampleblank

\begin{example}
设 $f(x)$ 在 $\lbrack 0,1\rbrack$ 上二次可导,且 $f(0)=0$,则存在 $\xi\in(0,1)$,使得
\[
    f'(\xi)-(\xi-1)^2f''(\xi)=0.
\]
\end{example}
\exampleblank

\begin{example}
设 $f(x),g(x)$ 在 $\lbrack a,b\rbrack$ 二次可导,且 $g''(x)\neq0$,若有
\[
    f(a)=f(b)=g(a)=g(b)=0,
\]
则存在 $\xi\in(a,b)$,使得
\[
    \dfrac{f(\xi)}{g(\xi)}=\dfrac{f''(\xi)}{g''(\xi)}.
\]
\end{example}
\exampleblank

\begin{example}
设 $f(x),g(x)$ 在 $\lbrack a,b\rbrack$ 上可导,且在 $(a,b)$ 上 $g'(x)\neq0$,则存在 $\xi\in(a,b)$,使得
\[
    \dfrac{f(a)-f(\xi)}{g(\xi)-g(b)}=\dfrac{f'(\xi)}{g'(\xi)}.
\]
\end{example}
\exampleblank

\begin{example}
设 $f\in C(\lbrack 0,1\rbrack)$,在 $(0,1)$ 上可导且 $f(1)=0$,则存在 $\xi\in(0,1)$ 使得
\[
    f'(\xi)=\left(1-\dfrac{1}{\xi}\right)f(\xi).
\]
\end{example}
\exampleblank

\begin{example}
设 $f\in C(\lbrack a,b\rbrack)$,在 $(a,b)$ 上可导,若有
\[
    f(a)f(b)>0,
    \qquad
    f(a)f\left(\dfrac{a+b}{2}\right)<0,
\]
则对任给的 $\lambda\in(-\infty,+\infty)$,存在 $\xi\in(a,b)$,使得
\[
    f'(\xi)=\lambda f(\xi).
\]
\end{example}
\exampleblank

\begin{example}
设 $f(x)$ 在 $\lbrack 0,+\infty)$ 上可导,若
\[
    f^2(a)-f^2(b)=b^2-a^2,
\]
则存在 $\xi\in(0,+\infty)$,使得
\[
    f'(\xi)=\dfrac{1-\xi^2}{(1+\xi^2)^2}.
\]
\end{example}
\exampleblank

\begin{example}
设 $f(x)$ 在 $\lbrack 0,1\rbrack$ 上二次可导,且 $f(0)=f(1)=0$,则存在 $\xi\in(0,1)$,使得
\[
    f''(\xi)=\dfrac{2f'(\xi)}{1-\xi}.
\]
\end{example}
\exampleblank

\begin{example}
设 $f(x)$ 在 $\lbrack 0,1\rbrack$ 上可导,且 $f(0)=0,f(x)>0\ (0<x<1)$,则存在 $\xi\in(0,1)$,使得
\[
    \dfrac{2f'(\xi)}{f(\xi)}
    =\dfrac{f'(1-\xi)}{f(1-\xi)}.
\]
\end{example}
\exampleblank


\newpage

\subsection{待定常数法(中值问题)}

该方法的本质是将需要证的中值设为常数,而后将原式子中的某个常数设为变量,再根据 Rolle 或 Lagrange 中值定理证明中值点的存在性.

\begin{example}
设 $f(x)$ 在 $\lbrack a,b\rbrack$ 上二次可导,则存在 $\xi\in(a,b)$,使得
\[
    f(b)-2f\left(\dfrac{a+b}{2}\right)+f(a)
    =\dfrac{(b-a)^2}{4}f''(\xi).
\]
\end{example}
\exampleblank

\begin{example}
设 $f(x),g(x)$ 在 $\lbrack a,b\rbrack$ 上二次可导,则存在 $\xi\in(a,b)$,使得
\[
    \dfrac{f(b)-f(a)-(b-a)f'(a)}{g(b)-g(a)-(b-a)g'(a)}
    =\dfrac{f''(\xi)}{g''(\xi)}.
\]
\end{example}
\exampleblank

\begin{example}
若存在 $f''(0)$,则
\[
    \lim\limits_{h\to0}
    \dfrac{f(2h)-2f(0)+f(-2h)}{4h^2}
    =f''(0).
\]
\end{example}
\exampleblank

\begin{example}
设 $h>0$,$f\in C(\lbrack a-h,a+h\rbrack)$,在 $(a-h,a+h)$ 上可导,则存在 $\theta:0<\theta<1$,使得
\[
    \dfrac{f(a+h)-f(a-h)}{h}
    =f'(a+\theta h)+f'(a-\theta h).
\]
\end{example}
\exampleblank

\begin{example}
设 $f(x)$ 在 $\lbrack a,b\rbrack$ 上连续,且在 $(a,b)$ 上二次可导,则对 $x\in(a,b)$,存在 $\xi\in(a,b)$,使得
\[
    \dfrac{f(x)-f(a)}{x-a}
    -\dfrac{f(b)-f(a)}{b-a}
    =\dfrac{1}{2}(x-b)f''(\xi).
\]
\end{example}
\exampleblank

\begin{example}
设 $f(x)$ 在 $\lbrack a,c\rbrack$ 上二次可导,$a<b<c$,则存在 $\xi\in(a,c)$,使得
\[
    \dfrac{f(a)}{(a-b)(a-c)}
    +\dfrac{f(b)}{(b-c)(b-a)}
    +\dfrac{f(c)}{(c-a)(c-b)}
    =\dfrac{1}{2}f''(\xi).
\]
\end{example}
\exampleblank

\begin{example}
设 $f(x)$ 在 $(a,b)$ 上三次可导,则存在 $\xi\in(a,b)$,使得
\[
    f(a)-f(b)+\dfrac{1}{2}(b-a)[f'(a)+f'(b)]
    =f'''(\xi)\dfrac{(b-a)^3}{12}.
\]
\end{example}
\exampleblank


\newpage

\subsection{利用几何意义(中值问题)}

Lagrange 中值定理的几何意义即,曲线上任意两点的弦,必与两点之间某点的切线平行.

\begin{example}
设 $f(x)$ 是可微函数,导函数 $f'(x)$ 严格单调递增,若 $f(a)=f(b)\ (a<b)$,试证明:对一切 $x\in(a,b)$,有
\[
    f(x)<f(a)=f(b).
\]
\end{example}
\exampleblank

\begin{example}
设 $f(x)$ 在区间 $\lbrack 0,1\rbrack$ 上可微,$f(0)=0,f(1)=1$,$k_1,k_2,\ldots,k_n$ 为 $n$ 个正数,证明在区间 $\lbrack 0,1\rbrack$ 上存在一组互不相等的数 $x_1,x_2,\ldots,x_n$,使得
\[
    \sum\limits_{i=1}^n\dfrac{k_i}{f'(x_i)}
    =\sum\limits_{i=1}^n k_i.
\]
\end{example}
\exampleblank

\begin{example}
设函数 $f(x)$ 在闭区间 $\lbrack a,b\rbrack$ 上连续,在开区间 $(a,b)$ 内可导,又 $f(x)$ 不是线性函数,且 $f(b)>f(a)$,试证:存在 $\xi\in(a,b)$,使得
\[
    f'(\xi)>\dfrac{f(b)-f(a)}{b-a}.
\]
\end{example}
\exampleblank

\begin{example}
设 $f(x)$ 在 $\lbrack a,b\rbrack$ 上连续,在 $(a,b)$ 内可导,$f(x)$ 既不是常数也不是一次函数,试证:存在 $\xi\in(a,b)$,使得
\[
    |f'(\xi)|>\left|\dfrac{f(b)-f(a)}{b-a}\right|.
\]
\end{example}
\exampleblank


\newpage

\subsection{多中值问题}

对于多个中值点的问题显然前面的常规方法失效了,此时要用到的是 Cauchy 中值定理:
\[
    F(b)-F(a)=\dfrac{F'(\xi)}{G'(\xi)}[G(b)-G(a)].
\]
若选取不同的函数 $G$,便可把 $F(b)-F(a)$ 表示成不同的形式.如另取 $G_1(x)$,则存在 $\xi,\eta\in(a,b)$,使得
\[
    \dfrac{F'(\xi)}{G'(\xi)}[G(b)-G(a)]
    =\dfrac{F'(\eta)}{G_1'(\eta)}[G_1(b)-G_1(a)].
\]
一般地,当 $G(x)$ 取 $n$ 个不同的函数时,便可得到含 $n$ 个中值的 $n-1$ 个等式.

\begin{example}
设 $f(x)$ 在 $\lbrack a,b\rbrack$ 上连续,在 $(a,b)$ 内可导,$f(a)\neq f(b)$,证明:存在 $\xi,\eta\in(a,b)$ 使得
\[
    f'(\xi)=\dfrac{a+b}{2\eta}f'(\eta).
\]
\end{example}
\exampleblank

\begin{example}
设 $f(x)$ 在 $\lbrack a,b\rbrack\ (b>a>0)$ 上连续,在 $(a,b)$ 内可导,证明:在 $(a,b)$ 内存在两点 $\xi,\eta$,使得
\[
    f'(\xi)=\dfrac{\eta^2}{ab}f'(\eta).
\]
\end{example}
\exampleblank

\begin{example}
设 $f(x)$ 在 $\lbrack a,b\rbrack$ 上连续,在 $(a,b)$ 内可导,$0<a<b$,证明:在 $(a,b)$ 内存在 $x_1,x_2,x_3$,使得
\[
    \dfrac{f'(x_1)}{2x_1}
    =(b^2+a^2)\dfrac{f'(x_2)}{4x_3^2}
    =\dfrac{\ln\dfrac{b}{a}}{b^2-a^2}x_3f'(x_3).
\]
\end{example}
\exampleblank

\begin{example}
设 $f(x)$ 在区间 $\lbrack a,b\rbrack$ 上连续,在 $(a,b)$ 内可导,且 $a\geq0$(或 $b\leq0$),试证:
\begin{enumerate}[label=(\arabic*)]
    \item 存在 $x_1,x_2,x_3\in(a,b)$,使得
    \[
        f'(x_1)
        =(b+a)\dfrac{f'(x_2)}{2x_2}
        =(b^2+ba+a^2)\dfrac{f'(x_3)}{3x_3^2};
    \]
    \item 对任意 $n\in\{1,2,\ldots\}$,存在 $x_1,x_2,\ldots,x_n\in(a,b)$,使得
    \[
    \begin{aligned}
        f'(x_1)
        &=(b+a)\dfrac{f'(x_2)}{2x_2}
        =\cdots\\
        &=(b^{n-1}+b^{n-2}a+\cdots+ba^{n-2}+a^{n-1})
        \dfrac{f'(x_n)}{nx_n^{n-1}}.
    \end{aligned}
    \]
\end{enumerate}
\end{example}
\exampleblank

\begin{example}
设 $f\in C(\lbrack a,b\rbrack)$,且在 $(a,b)$ 上可导,若 $f(a)=f(b)=1$,则存在 $\xi,\eta\in(a,b)$,使得
\[
    e^{\eta-\xi}[f(\eta)+f'(\eta)]=1.
\]
\end{example}
\exampleblank

\begin{example}
设 $f\in C(\lbrack a,b\rbrack)$,且在 $(a,b)$ 上可导,若 $f(0)=0,f(1)=1$,则存在 $\xi,\eta\in(a,b)$,使得
\[
    f'(\xi)f'(\eta)=1,
    \qquad \xi\neq\eta.
\]
\end{example}
\exampleblank


\newpage

\subsection{函数与导函数的关系}

设 $f(x)$ 在 $(a,+\infty)$ 上可导,如果只有极限 $\displaystyle\lim\limits_{x\to+\infty}f'(x)=B$ 存在,那么当 $x\to+\infty$ 时,$f(x)$ 不一定存在极限,例如 $f(x)=\sin(\ln x)$;反之,即使存在 $\displaystyle\lim\limits_{x\to+\infty}f(x)$,也不一定存在 $\displaystyle\lim\limits_{x\to+\infty}f'(x)$,例如 $f(x)=\dfrac{\sin x^2}{x}$.但若二者都存在,那么必有 $\displaystyle\lim\limits_{x\to+\infty}f'(x)=0$.

\begin{example}
设 $f(x)$ 在 $(a,+\infty)$ 上可导,若 $\displaystyle\lim\limits_{x\to+\infty}f'(x)=B>0$,则:
\begin{enumerate}[label=(\arabic*)]
    \item $\displaystyle\lim\limits_{x\to+\infty}f(x)=+\infty$;
    \item $f(x)\sim Bx\ (x\to+\infty)$.
\end{enumerate}
\end{example}
\exampleblank

\begin{example}[\markstar]
设 $f\in C^{(1)}((a,+\infty))$,若 $f'(x)$ 在 $(a,+\infty)$ 上一致连续,且存在 $\displaystyle\lim\limits_{x\to+\infty}f(x)$,则
\[
    \lim\limits_{x\to+\infty}f'(x)=0.
\]
\end{example}
\exampleblank

\begin{example}[\markstar]
设 $f(x)$ 在 $\lbrack 0,+\infty)$ 上三阶可导,且有
\[
    \lim\limits_{x\to+\infty}f(x)=A,
    \qquad
    \lim\limits_{x\to+\infty}f'''(x)=0,
\]
则
\[
    \lim\limits_{x\to+\infty}f'(x)
    =\lim\limits_{x\to+\infty}f''(x)=0.
\]
\end{example}
\exampleblank

\begin{example}[\markstar]
试证明下列命题:
\begin{enumerate}[label=(\arabic*)]
    \item 设 $f(x)$ 在 $(0,+\infty)$ 上可导,$a>0$,若
    \[
        \lim\limits_{x\to+\infty}[af(x)+f'(x)]=l,
    \]
    则
    \[
        \lim\limits_{x\to+\infty}f(x)=\dfrac{l}{a};
    \]
    \item 设 $f(x)$ 在 $(0,+\infty)$ 上可导,$a>0$,若有
    \[
        \lim\limits_{x\to+\infty}[af(x)+2\sqrt{x}f'(x)]=l,
    \]
    则
    \[
        \lim\limits_{x\to+\infty}f(x)=\dfrac{l}{a};
    \]
    \item 设 $f(x)$ 在 $(0,+\infty)$ 上二次可导,若有
    \[
        \lim\limits_{x\to+\infty}[f(x)+f'(x)+f''(x)]=l,
    \]
    则
    \[
        \lim\limits_{x\to+\infty}f(x)=l.
    \]
\end{enumerate}
\end{example}
\exampleblank

\begin{example}
设 $f(x)$ 在 $(0,+\infty)$ 上可导,若 $\displaystyle\lim\limits_{x\to+\infty}f'(x)=l$,则对任意的 $A$ 值,有
\[
    \lim\limits_{x\to+\infty}[f(x+A)-f(x)]=lA.
\]
\end{example}
\exampleblank

\begin{example}[\markstar]
设 $f(x)$ 在 $\lbrack 0,+\infty)$ 上可微,$f(0)=0$,且
\[
    |f'(x)|\leq A|f(x)|,
    \qquad x\in\lbrack 0,+\infty),
\]
则 $f\equiv0$.
\end{example}
\exampleblank

\begin{example}
设 $f(x),g(x)$ 在 $\lbrack a,b\rbrack$ 上连续,$g(x)$ 在 $(a,b)$ 内可微,$g(a)=0$,若有实数 $\lambda\neq0$,使得
\[
    |g(x)f(x)+\lambda g'(x)|\leq|g(x)|,
    \qquad x\in(a,b),
\]
成立,则 $g\equiv0$.
\end{example}
\exampleblank

\begin{example}
证明:在 $\mathbb{R}$ 上的任何可微函数都不可能满足下列函数方程:
\begin{enumerate}[label=(\arabic*)]
    \item $g(g(x))=-x^3+x+1$;
    \item $g(g(x))=-x^3+x^2\pm1$;
    \item $g(g(x))=x^2-3x+3$.
\end{enumerate}
\end{example}
\exampleblank


\newpage

\section{Taylor 公式}


\subsection{中值等式}

\begin{example}
设 $f(x)$ 在 $\lbrack a,b\rbrack$ 上三次可导,试证:存在 $c\in(a,b)$,使得
\[
    f(b)=f(a)+f'\left(\dfrac{a+b}{2}\right)(b-a)
    +\dfrac{1}{24}f'''(c)(b-a)^3.
\]
\end{example}
\exampleblank

\begin{example}
设 $f\in C^{(3)}(\lbrack -1,1\rbrack)$,且 $f(-1)=0,f(1)=1,f'(0)=0$,则存在 $\xi\in(-1,1)$,使得
\[
    f'''(\xi)=3.
\]
\end{example}
\exampleblank

\begin{example}
设 $f(x)$ 在 $\lbrack -1,1\rbrack$ 上二次可导,若有 $f(-1)=0,f(0)=0,f(1)=1$,则存在 $\xi\in(-1,1)$,使得
\[
    f''(\xi)=1.
\]
\end{example}
\exampleblank

\begin{example}
设 $f(x)$ 在 $\lbrack -1,1\rbrack$ 上三次可导,若有 $f(-1)=f(0)=0,f(1)=1,f'(0)=0$,则存在 $\xi\in(-1,1)$,使得
\[
    f'''(\xi)\geq 3.
\]
\end{example}
\exampleblank


\newpage

\subsection{界的估计}

\begin{example}
设 $f(x)$ 二次可微,且 $f(0)=f(1)=0$,又
$\displaystyle\max\limits_{0\leq x\leq1}f(x)=2$,试证:
\[
    \min\limits_{0\leq x\leq1}f''(x)\leq-16.
\]
\end{example}
\exampleblank

\begin{example}
已知 $f(x)$ 在 $\lbrack 0,1\rbrack$ 上二阶可导,且 $f(0)=0,f(1)=3$,又
$\displaystyle\min\limits_{x\in\lbrack 0,1\rbrack}f(x)=-1$,证明:存在 $c\in(0,1)$,使得
\[
    f''(c)\geq18.
\]
\end{example}
\exampleblank

\begin{example}[\markstar]
设 $f(x)$ 在 $\lbrack 0,1\rbrack$ 上二次可导,若有
\[
    |f(x)|\leq A,\qquad |f''(x)|\leq B,\qquad x\in\lbrack 0,1\rbrack,
\]
则
\[
    |f'(x)|\leq2A+\dfrac{B}{2}.
\]
\end{example}
\exampleblank

\begin{example}
设 $f(x)$ 在 $\lbrack 0,1\rbrack$ 上二阶可导,且有 $f(0)=f(1)$,$|f''(x)|\leq M$,则
\[
    |f'(x)|\leq\dfrac{M}{2}.
\]
\end{example}
\exampleblank

\begin{example}[\markstar]
设 $f(x)$ 在 $(-\infty,+\infty)$ 上二次可导,且有
\[
    M_k=\sup\{|f^{(k)}(x)|:-\infty<x<+\infty\}<+\infty.
\]
则
\[
    M_1\leq\sqrt{2M_0M_2}.
\]
\end{example}
\exampleblank

\begin{example}[\markstar]
设 $f(x)$ 在 $(-\infty,+\infty)$ 上 $m$ 次可导,且有
\[
    M_k=\sup\{|f^{(k)}(x)|:-\infty<x<+\infty\}<+\infty.
\]
则
\[
    M_k\leq2^{\frac{k(m-k)}{2}}M_0^{1-\frac{k}{m}}M_m^{\frac{k}{m}},
    \qquad k=1,2,\ldots,m-1.
\]
\end{example}
\exampleblank


\newpage

\subsection{中值点的极限}

\begin{example}
证明:对每个 $x>0$,均存在 $\theta(x)\in\lbrack \dfrac{1}{4},\dfrac{1}{2}\rbrack$,使得
\[
    \sqrt{x+1}-\sqrt{x}=\dfrac{1}{2\sqrt{x+\theta(x)}}.
\]
且
\[
    \lim\limits_{x\to0^+}\theta(x)=\dfrac{1}{4},
    \qquad
    \lim\limits_{x\to+\infty}\theta(x)=\dfrac{1}{2}.
\]
\end{example}
\exampleblank

\begin{example}
设 $f(x)$ 在点 $a$ 的某个邻域上 $n+2$ 阶连续可微,且 $f^{(n+2)}(a)\neq0$,并有
\[
    f(a+h)=f(a)+f'(a)h+\cdots+\dfrac{f^{(n)}(a)}{n!}h^n
    +\dfrac{f^{(n+1)}(a+\theta h)}{(n+1)!}h^{n+1},
\]
其中 $h\in(0,\delta)$,$0<\theta<1$.证明:
\[
    \lim\limits_{h\to0}\theta=\dfrac{1}{n+2}.
\]
\end{example}
\exampleblank

\begin{example}
设 $f(x)$ 在 $(x_0-\delta,x_0+\delta)$ 上 $n$ 阶连续可微,已知
\[
    f^{(2)}(x_0)=f^{(3)}(x_0)=\cdots=f^{(n-1)}(x_0)=0,
    \qquad f^{(n)}(x_0)\neq0,
\]
若 $0<|h|<\delta$ 时有
\[
    \dfrac{f(x_0+h)-f(x_0)}{h}=f'(x_0+\theta h),
\]
证明:
\[
    \lim\limits_{h\to0}\theta
    =\sqrt[n-1]{\dfrac{1}{n}}.
\]
\end{example}
\exampleblank

\begin{example}
设 $f(x)$ 在点 $a$ 处二阶可导,且 $f''(a)=0$.证明当 $h$ 充分小时,有等式
\[
    f(a+h)-f(a)=f'(a+\theta h)h,
\]
其中 $\theta\in(0,1)$,并且
\[
    \lim\limits_{h\to0}\theta=\dfrac{1}{2}.
\]
\end{example}
\exampleblank


\newpage

\subsection{极值问题}

设 $f(x)$ 在 $x_0$ 处二阶可导,且 $f'(x_0)=0$.若 $f''(x_0)<0$($>0$),则 $x_0$ 为 $f(x)$ 的极大(小)值点.

\begin{example}
设 $f\in C(\lbrack a,b\rbrack)$,且 $f\in C^{(2)}((a,b))$,若有
\[
    f''(x)=e^x f(x),\qquad f(a)=f(b)=0,
\]
则 $f(x)\equiv0$.
\end{example}
\exampleblank

\begin{example}
设 $f\in C^{(2)}((-\infty,+\infty))$,若有
\[
    [f(0)]^2+[f'(0)]^2=4,
    \qquad |f(x)|=1\quad(-\infty<x<+\infty),
\]
则存在 $\xi$,使得
\[
    f(\xi)+f''(\xi)=0.
\]
\end{example}
\exampleblank

\begin{example}
设 $f(x),g(x)$ 定义在 $(-\infty,+\infty)$ 上,且 $f(x)$ 二次可导,又有
\[
    f(a)=f(b)=0\quad(a<b),
\]
且
\[
    f''(x)+f'(x)g(x)-f(x)=0,
    \qquad x\in(-\infty,+\infty),
\]
则 $f(x)\equiv0$.
\end{example}
\exampleblank

\begin{example}
设 $f\in C^{(2)}(\lbrack 0,a\rbrack)$,且 $x_0\in(0,a)$ 是 $f(x)$ 的极小值点,则
\[
    |f'(0)|+|f'(a)|
    \leq a\max\limits_{x\in\lbrack 0,a\rbrack}|f''(x)|.
\]
\end{example}
\exampleblank


\newpage

\section{凹凸性与不等式}


\subsection{凹凸性}

\begin{definition}
设 $f(x)$ 在区间 $I$ 上有定义,$f(x)$ 在 $I$ 上称为凸函数,当且仅当对任意 $x_1,x_2\in I$ 以及任意 $\lambda\in(0,1)$,有
\[
    f(\lambda x_1+(1-\lambda)x_2)
    \leq\lambda f(x_1)+(1-\lambda)f(x_2).
\]
\end{definition}

\begin{definition}
设 $f(x)$ 在区间 $I$ 上有定义,$f(x)$ 在 $I$ 上称为凸函数,当且仅当对任意 $x_1,x_2\in I$,有
\[
    f\left(\dfrac{x_1+x_2}{2}\right)
    \leq\dfrac{f(x_1)+f(x_2)}{2}.
\]
\end{definition}

\begin{definition}
设 $f(x)$ 在区间 $I$ 上有定义,$f(x)$ 在 $I$ 上称为凸函数,当且仅当对任意 $x_1,x_2,\ldots,x_n\in I$,有
\[
    f\left(\dfrac{x_1+x_2+\cdots+x_n}{n}\right)
    \leq\dfrac{f(x_1)+f(x_2)+\cdots+f(x_n)}{n}.
\]
\end{definition}

\begin{definition}
设 $f(x)$ 在区间 $I$ 上有定义,当且仅当曲线 $y=f(x)$ 的切线保持在曲线以下,称 $f(x)$ 为凸函数.
\end{definition}

\begin{lemma}[\markstar]
设 $f(x)$ 是区间 $I$ 上的凸函数,则对 $I$ 的任一内点 $x$,单侧导数 $f'_+(x),f'_-(x)$ 皆存在,皆为增函数,且
\[
    f'_-(x)\leq f'_+(x),
    \qquad x\in\mathring I.
\]
\end{lemma}

因此,更一般地,有如下定理.

\begin{theorem}[\markstar]
$f$ 为凸函数,当且仅当对任意 $x_0\in\mathring I$,存在 $\alpha\in\mathbb{R}$,使得
\[
    f(x)\geq\alpha(x-x_0)+f(x_0).
\]
\end{theorem}

\begin{remark}
当 $f$ 可导时,显然可取 $\alpha=f'(x_0)$.
\end{remark}

\begin{theorem}[\markstar]
\begin{enumerate}[label=(\arabic*)]
    \item 一般情况下,定义 4.4.2 与定义 4.4.3 等价;
    \item $f$ 连续时,定义 4.4.1、定义 4.4.2 与定义 4.4.3 等价;
    \item $f$ 可导时,定义 4.4.1、定义 4.4.2、定义 4.4.3、定义 4.4.4 与 $f'$ 单增等价;
    \item $f$ 二阶可导时,定义 4.4.1、定义 4.4.2、定义 4.4.3、定义 4.4.4、$f'$ 单增与 $f''>0$ 等价.
\end{enumerate}
\end{theorem}

\begin{theorem}[\markstar]
设 $f(x)$ 为 $\lbrack a,b\rbrack$ 上的凸函数,则:
\begin{enumerate}[label=(\arabic*)]
    \item $f$ 在 $(a,b)$ 上连续;
    \item $f(x)\leq\max\{f(a),f(b)\}$.
\end{enumerate}
\end{theorem}

\begin{remark}
事实上,凸函数在定义域内部有良好性质,但边界上无法约束(就算端点特别大或特别小,亦不影响其凹凸性).
\end{remark}

\begin{example}
设 $f(x)$ 为区间 $(a,b)$ 内的凸函数,试证:$f(x)$ 在 $I$ 的任一闭区间 $\lbrack \alpha,\beta\rbrack\subseteq(a,b)$ 上满足 Lipschitz 条件.
\end{example}
\exampleblank

\begin{example}
设 $f(0)=0$,$f(x)$ 在 $\lbrack 0,+\infty)$ 上为非负的严格凸函数,令
\[
    F(x)=\dfrac{f(x)}{x},\qquad x>0,
\]
证明:$f(x),F(x)$ 为严格递增的.
\end{example}
\exampleblank

\begin{example}
设 $f(x),g(x)$ 在区间 $(a,b)$ 内有定义,且对任意 $x,x_0\in(a,b)$ 有
\[
    f(x)-f(x_0)\geq g(x_0)(x-x_0).
\]
试证:对任意 $x_0\in(a,b)$,$f(x)$ 在 $x_0$ 处连续,并且有左右单侧导数.
\end{example}
\exampleblank

\begin{example}
证明下列命题成立:
\begin{enumerate}[label=(\arabic*)]
    \item 两凸函数之和仍为凸函数;
    \item 两递增非负凸函数之积仍为凸函数.
\end{enumerate}
\end{example}
\exampleblank

\begin{example}
设 $f\in C(\lbrack 0,1\rbrack)$,且在 $(0,1)$ 上二次可导,若有
\[
    f(0)=f(1)=0,
    \qquad f''(x)+2f'(x)+f(x)=0\quad(0<x<1),
\]
则 $f(x)\leq0$.
\end{example}
\exampleblank

\begin{example}
设 $f(x)$ 在 $(a,b)$ 上四次可导,$x_0\in(a,b)$,且有
\[
    f^{(2)}(x_0)=0,
    \qquad f^{(3)}(x_0)=0,
    \qquad f^{(4)}(x_0)>0,
\]
则 $f(x)$ 是下凸函数.
\end{example}
\exampleblank

\begin{example}
设 $f(x)$ 在 $(-\infty,+\infty)$ 上二次可导,且有
\[
    f(x)\leq0,
    \qquad f''(x)\geq0,
\]
则 $f(x)\equiv C$(常数).
\end{example}
\exampleblank

\begin{example}
设 $f(x)$ 在 $(a,b)$ 上可导,若对 $(a,b)$ 中的 $x,y$($x\neq y$),存在唯一的 $\xi\in(a,b)$,使得
\[
    \dfrac{f(y)-f(x)}{y-x}=f'(\xi),
\]
则 $f(x)$ 是严格上凸或下凸的.
\end{example}
\exampleblank

\begin{example}
设 $f(x)$ 是 $(0,1)$ 上三次可导的非负函数,若存在 $x_1,x_2\in(0,1)$ 使得 $f(x_1)=f(x_2)=0$,则存在点 $\xi\in(0,1)$,使得
\[
    f'''(\xi)=0.
\]
\end{example}
\exampleblank


\newpage

\subsection{不等式}

基础配套教材已经做了基本方法总结,接下来以实例来讲解具体如何运用.


\newpage

\methodsection{(A) 利用单调性}

\begin{proposition}[\marktwostar]
当 $x\in\left(0,\dfrac{\pi}{2}\right)$ 时,
\[
    \dfrac{2}{\pi}x<\sin x<x.
\]
\end{proposition}

\begin{example}
对任意 $n=1,2,3,\ldots$,求能使不等式
\[
    \left(1+\dfrac{1}{n}\right)^{n+\alpha}
    \leq e
    \leq\left(1+\dfrac{1}{n}\right)^{n+\beta}
\]
成立的 $\alpha$ 的最大值与 $\beta$ 的最小值.
\end{example}
\exampleblank

\begin{example}
证明如下不等式:
\begin{enumerate}[label=(\arabic*)]
    \item 当 $0<x_1<x_2<\dfrac{\pi}{2}$ 时,有
    \[
        \dfrac{\tan x_2}{\tan x_1}>\dfrac{x_2}{x_1};
    \]
    \item 当 $x\in\left(0,\dfrac{\pi}{2}\right)$ 时,有
    \[
        \dfrac{\tan x}{x}>\dfrac{x}{\sin x};
    \]
    \item 当 $x\in\left(0,\dfrac{\pi}{2}\right)$ 时,有
    \[
        2\sin x+3\tan x>3x.
    \]
\end{enumerate}
\end{example}
\exampleblank

\begin{example}
\begin{enumerate}[label=(\arabic*)]
    \item 设 $b>a>e$,证明 $a^b>b^a$;
    \item 比较 $\pi^e$ 与 $e^\pi$ 的大小.
\end{enumerate}
\end{example}
\exampleblank


\newpage

\methodsection{(B) 利用中值定理}

\begin{example}
证明:
\begin{enumerate}[label=(\arabic*)]
    \item 当 $x>0$ 时,
    \[
        \ln\left(1+\dfrac{1}{x}\right)>\dfrac{1}{x+1};
    \]
    \item 当 $x>0$ 时,
    \[
        \dfrac{x}{1+x^2}<\arctan x<x.
    \]
\end{enumerate}
\end{example}
\exampleblank

\begin{example}
\begin{enumerate}[label=(\arabic*)]
    \item 证明:当 $s>0$ 时,
    \[
        \dfrac{n^{s+1}}{s+1}
        <1^s+2^s+\cdots+n^s
        <\dfrac{(n+1)^{s+1}}{s+1};
    \]
    \item 设 $\displaystyle\lambda=\sum\limits_{k=1}^n\dfrac{1}{k}$,证明
    \[
        e^\lambda>n+1.
    \]
\end{enumerate}
\end{example}
\exampleblank


\newpage

\methodsection{(C) 利用凹凸性}

\begin{example}
证明下列不等式:
\begin{enumerate}[label=(\arabic*)]
    \item 当 $x\in\lbrack 0,1\rbrack$ 时,
    \[
        1+x^2\leq2^x\leq1+x;
    \]
    \item 当 $x\in\lbrack 0,1\rbrack$ 时,
    \[
        \sin(\pi x)\leq\dfrac{\pi^2}{2}x(1-x).
    \]
\end{enumerate}
\end{example}
\exampleblank

\begin{example}[\markstar]
证明下列经典不等式:
\begin{enumerate}[label=(\arabic*)]
    \item (平均值不等式)设 $a_i>0$,则
    \[
        \dfrac{n}{\displaystyle\sum\limits_{i=1}^n\dfrac{1}{a_i}}
        \leq\sqrt[n]{a_1a_2\cdots a_n}
        \leq\dfrac{\displaystyle\sum\limits_{i=1}^n a_i}{n}.
    \]
    \item (Young 不等式)设 $a,b>0$,$0<p,q<1$,且 $\dfrac{1}{p}+\dfrac{1}{q}=1$,则
    \[
        ab\leq\dfrac{a^p}{p}+\dfrac{b^q}{q}.
    \]
    \item (Holder 不等式)设 $a_i,b_i>0$,$0<p,q<1$,且 $\dfrac{1}{p}+\dfrac{1}{q}=1$,则
    \[
        \sum\limits_{i=1}^n a_i b_i
        \leq
        \left(\sum\limits_{i=1}^n a_i^p\right)^{\frac{1}{p}}
        \left(\sum\limits_{i=1}^n b_i^q\right)^{\frac{1}{q}}.
    \]
\end{enumerate}
\end{example}
\exampleblank

\end{document}
